\documentclass[exercice,a4]{thedocument}%
\startdatakeys
    \source{laclassecooperativedestef.com}
    \BO{2026}
    \cycle{3}
    \competencesDuSocle{RE4,CO2}
    \objectifsApprentissage{C3C1ab,C3C1aj,C3C1ad,C3C2ac,C3C2ab,C3C1ae,C3C1ag}
    \motsCles{milieu, cercle, perpendicularite, quadrilatere}
\enddatakeys
\begin{document}%
\noindent\textbf{Programme de construction 22}\par\smallskip\noindent
\begin{itemize}
    \item Tracer la droite $(\mathspoint{d})$.
    \item Placer $\mathspoint{A}$ et $\mathspoint{B}$ sur
    $(\mathspoint{d})$ tels que $\mathspoint{AB}=4\text{ cm}$.
    \item Tracer le cercle $\mathcal{A}$ de centre $\mathspoint{A}$ et de
    rayon $2\text{ cm}$.
    \item Tracer le cercle $\mathcal{B}$ de centre $\mathspoint{B}$ et de
    rayon $2\text{ cm}$.
    \item Placer $\mathspoint{C}$, milieu de $[\mathspoint{AB}]$.
    \item Tracer la droite $(\mathspoint{e})$ perpendiculaire à
    $(\mathspoint{d})$ passant par $\mathspoint{C}$.
    \item Placer les points $\mathspoint{D}$ et $\mathspoint{E}$ sachant
    que $\mathspoint{DC}=2\text{ cm}$ et $\mathspoint{CE}=3\text{ cm}$.
    \item Tracer le quadrilatère $\mathspoint{EBDA}$.
\end{itemize}
\end{document}%
